%! Justifying a Claim Based on a Confidence Interval for a Population Proportion — Unit 6, Lesson 6.3 — Homework
\documentclass[10pt]{article}
\usepackage{apstats-article}
\usepackage{apstats-boxes}
\newcommand{\TallMath}[1]{$\displaystyle #1\rule[-1.4em]{0pt}{3.2em}$}

\begin{document}

\pageheader{Unit 6, Lesson 6.3}{Homework}
\namedateperiod

\vspace{0.15in}

\begin{enumerate}

  \item A 95\% CI for the proportion of adults who exercise at least 3 days per week is
        $(0.42, 0.54)$.
        \begin{enumerate}[label=\alph*.]
          \item Write a complete interpretation of this CI. \writelines{2}
          \item A fitness organization claims 50\% of adults exercise at least 3 days per week.
                Is this plausible? \writelines{2}
          \item The surgeon general claims \emph{fewer than} 45\% of adults exercise regularly.
                Does the CI support this claim? \writelines{2}
        \end{enumerate}

  \item A 99\% CI for the proportion of students who own a calculator is $(0.71, 0.89)$.
        \begin{enumerate}[label=\alph*.]
          \item A teacher claims 90\% of students own a calculator. Plausible? \writelines{2}
          \item Would a 95\% CI using the same data be wider or narrower? Explain.
                \writelines{2}
          \item If the sample size doubled (holding $\hat p$ constant), how would the width of
                the 99\% CI change? \writelines{2}
        \end{enumerate}

  \item \textbf{Correcting errors.} For each student statement, identify and correct the error.
        \begin{enumerate}[label=\alph*.]
          \item ``There is a 95\% probability that the true proportion of left-handed students
                is between 0.09 and 0.17.'' \writelines{2}
          \item ``95\% of all students are left-handed according to our interval.'' \writelines{2}
        \end{enumerate}

\end{enumerate}

\begin{reflectionbox}
\textbf{Preview:} Lesson 6.4 introduces a new inference tool --- the \emph{significance test} for
a population proportion. Instead of estimating $p$ with an interval, we will ask: is the sample
evidence strong enough to reject a specific claimed value of $p$? Think about how this differs
from saying a value is ``not plausible.''
\end{reflectionbox}

\end{document}
